\documentclass[11pt]{article}
\usepackage[a4paper,margin=30mm]{geometry}
\usepackage{amsmath,amssymb,amsthm}
\usepackage{booktabs}
\usepackage[hidelinks]{hyperref}
\usepackage{microtype}
\newtheorem{theorem}{Theorem}
\newtheorem{proposition}[theorem]{Proposition}
\newtheorem{lemma}[theorem]{Lemma}
\newtheorem{corollary}[theorem]{Corollary}
\newtheorem{remark}[theorem]{Remark}
\newcommand{\Q}{\mathbf Q}
\newcommand{\F}{\mathbf F}
\title{Square-class patterns of seven consecutive integers\\and three elementary integral quartics}
\author{\textbf{AUTHOR IDENTITY CONFIRMATION REQUIRED}}
\date{3 September 2026}
\begin{document}
\maketitle
\begin{abstract}
We study the affine rank of the square classes of seven consecutive nonzero
integers. This records when a common rational scaling makes every term a square
in a field of degree at most four generated by two square roots. After affine
rank at most one is excluded, enumeration up to relabelling gives 651 three-
or four-block equality patterns; quotienting by reversal and applying known
six-term and four-square restrictions leaves 284 candidates, each controlled
by precisely fifteen character quotients. A four-consecutive-factor identity
eliminates 186 patterns. We determine all integral points on
$y^2=t(t+2)(t+3)(t+6)$ by an elementary squarefree-kernel argument and finite
congruence certificates. Together with an integral reflection, this eliminates
44 further patterns after a same-parameter check, leaving 54. Two
constant-difference factorizations exclude another 19 and 12 patterns, leaving
23 explicitly listed necessary candidates. Every computational count and
congruence certificate is bound by a versioned reproducibility manifest. No
bounded search or unproved Mordell--Weil computation enters the proof. The
theorem neither proves that any remaining pattern is realizable nor decides the
unrestricted rational-progression problem.
\end{abstract}

\noindent\textbf{Keywords:} square classes; arithmetic progressions; integral
points; biquadratic fields; Diophantine equations.

\medskip
\noindent\textbf{Mathematics Subject Classification (2020):} Primary 11B25;
Secondary 11D45, 11G05.

\section{Square-class formulation}
Put $G=\Q^*/\Q^{*2}$, written additively over $\F_2$. For the nonzero
numbers $a_i=t+i$, define
\[
 \operatorname{arank}(a_0,\ldots,a_6)
 =\dim_{\F_2}\langle [a_i]-[a_0]:1\leq i\leq6\rangle .
\]
The parameters $t=-6,\ldots,0$ are excluded.

For reference, define
\[
 R_s(N)=\max_{\substack{a,q\in\Q,\ q\ne0;\ V\leq G,\ \dim_{\F_2}V=s\\
                         a+qi\ne0\ (0\leq i<N)}}
             \#\{0\leq i<N:[a+qi]\in V\},
\]
so every squareclass occurring in the definition is defined.  Thus the
maximum ranges over nonconstant rational arithmetic progressions whose
$N$ displayed terms are all nonzero, and over exactly $s$-dimensional
$\F_2$-subspaces of $G$.  The maximum is unchanged if one permits a common
rational scaling and replaces linear rank by affine rank.

\begin{lemma}\label{lem:R16}
$R_1(6)=5$.
\end{lemma}
\begin{proof}
If six rational terms have classes in $V=\langle[D]\rangle$, they are six
squares in the quadratic field $\Q(\sqrt D)$.  Xarles proves that a
nonconstant progression over a quadratic field has at most five square
terms \cite[Theorem~1, $S(2)=6$]{XarlesNF}.  Conversely,
\[
49,169,289,409,529
\]
is a nonconstant progression of squares over $\Q(\sqrt{409})$; appending
the sixth term $649$ shows that five of six terms meet the condition, so
the upper bound is attained.
\end{proof}

\begin{proposition}[Common scaling]\label{prop:scaling}
The seven classes have affine rank at most two if and only if there are a
common $\lambda\in\Q^*$ and a field
$L=\Q(\sqrt{D_1},\sqrt{D_2})$ of degree at most four such that every
$\lambda(t+i)$ is a square in $L$. If the difference classes span a
two-dimensional space, $L$ may be chosen biquadratic.
\end{proposition}
\begin{proof}
Write $[t+i]=c+v_i$ with $v_i\in
V=\langle[D_1],[D_2]\rangle$, and choose $[\lambda]=c$.
Then $[\lambda(t+i)]=v_i$ dies in $L^*/L^{*2}$. Conversely,
\[
 \ker\bigl(\Q^*/\Q^{*2}\longrightarrow L^*/L^{*2}\bigr)
 =\langle[D_1],[D_2]\rangle ,
\]
so the classes lie in a common affine coset of dimension at most two.
Indeed, if a rational class $[r]$ dies in $L$, then either $r$ is already
a rational square or $\Q(\sqrt r)$ is a quadratic subfield of $L$.
The quadratic subfields of a biquadratic extension are exactly
$\Q(\sqrt{D_1})$, $\Q(\sqrt{D_2})$, and
$\Q(\sqrt{D_1D_2})$, by the Kummer correspondence.  This gives the stated
kernel; if $[L:\Q]<4$, the same argument gives the corresponding
zero- or one-dimensional subspace.
\end{proof}
\begin{remark}
The scaling is essential: the common class $c$ can be independent of the
two difference classes, in which case adjoining $\sqrt c$ may raise the
degree to eight. We do not claim the unscaled $t+i$ are squares in one
fixed biquadratic field.
\end{remark}

\section{The finite pattern space}
Represent a rank-two configuration by labels $\ell_i\in\F_2^2$.
Lemma~\ref{lem:R16} forces both endpoint six-term windows to
contain at least three labels. A block of size at least five is excluded
by $Q(7)=4$, and a block at four positions $(i,j,k,l)$ with
$j-i=l-k$ is excluded by
\cite[Proposition~5 and Corollary~6]{GJXRudin}.

\begin{proposition}[Exact pattern count]\label{prop:count}
Before these screens there are 343 label patterns up to reversal. The
screens eliminate 59 reversal orbits, leaving 284.
\end{proposition}
\begin{proof}
Since $\operatorname{AGL}(2,2)\cong S_4$, an orbit using three or four
labels is an unlabelled partition into three or four blocks. Restricted
growth strings enumerate
\[
 S(7,3)+S(7,4)=301+350=651
\]
partitions. Reversal fixes 35, so Burnside gives $(651+35)/2=343$.

The finite screen discards a word if an endpoint six-window has at most
two blocks, if a block has size at least five, or if a four-element block
has equal first and last gaps. It discards 109 words forming 59 reversal
orbits. Of the remaining 542 words, 26 are reversal-fixed, giving
$(542+26)/2=284$. The supplementary generator directly asserts every
count.
\end{proof}

For a label word define
\[
 T:\F_2^7\longrightarrow\F_2\oplus\F_2^2,\qquad
 e\longmapsto\left(\sum e_i,\sum e_i\ell_i\right).
\]
\begin{lemma}[Fifteen characters]\label{lem:15}
For each pattern, $\ker T$ has dimension four and exactly 15 nonzero
elements. Every such $e$ gives the necessary equation
\[
 y_e^2=\prod_{e_i=1}(t+i).
\]
\end{lemma}
\begin{proof}
The labels affinely span $\F_2^2$, so $T$ has rank three. Rank--nullity
gives dimension four. If $[t+i]=c+\phi(\ell_i)$, the class of the product
is $(\sum e_i)c+\phi(\sum e_i\ell_i)=0$. These are all affine linear
relations.
\end{proof}

\section{The first 186 exclusions}
\begin{lemma}[Four consecutive factors]\label{lem:four}
If $x\in\mathbf Z$ and $x(x+1)(x+2)(x+3)$ is a rational square, then
$x\in\{-3,-2,-1,0\}$.
\end{lemma}
\begin{proof}
It is an integer square $y^2$. Put $n=x(x+3)$. Then
\[
 x(x+1)(x+2)(x+3)=(n+1)^2-1,
\]
so $(n+1-y)(n+1+y)=1$. The two integral factor pairs give $n=0$ or
$-2$, and the four asserted values follow.
\end{proof}
The consecutive masks are $15,30,60,120$. Recomputing all 15-element
kernels and intersecting with these masks eliminates exactly 186 of 284
patterns, leaving 98. This is finite exact computation, not height search.

\section{The mask 77 curve}
Consider
\[
 H:\quad y^2=t(t+2)(t+3)(t+6).
\]
It is birational to $E:Y^2=X^3-36X$ by
\[
 X=12+\frac{36}{t},\quad Y=\frac{36y}{t^2},\qquad
 t=\frac{36}{X-12},\quad y=\frac{36Y}{(X-12)^2}.
\]
The boundary correspondence is
\[
\begin{array}{c|c}
H&E\\ \hline
(0,0)&O\\ (-2,0)&(-6,0)\\ (-3,0)&(0,0)\\ (-6,0)&(6,0)\\
\infty_+&(12,36)\\ \infty_-&(12,-36).
\end{array}
\]
This is Cremona 576h2, equivalently LMFDB 576.c3
\cite{LMFDB576c3}. The label is only a cross-check: the condition
$36/(X-12)\in\mathbf Z$ is not ordinary integrality of $X$.

\begin{theorem}\label{thm:H}
\[
 H(\mathbf Z)=\{(-6,0),(-3,0),(-2,0),(0,0),(6,72),(6,-72)\}.
\]
\end{theorem}
\begin{proof}
Put $A=t(t+6)$ and $B=(t+2)(t+3)$. Since $B-A=6-t$ and
$A=(t+12)(t-6)+72$, one has $\gcd(A,B)\mid72$. At a nonbranch integral
point, $B>0$ and $AB=y^2>0$, hence $A>0$. Their common positive
squarefree kernel is
\[
 d\in\{1,2,3,6\},\qquad A=dU^2,\quad B=dV^2.
\]
With $x=t+3$ this becomes
\[
 x^2-dU^2=9,\qquad x(x-1)=dV^2.
\]
Write $d=ab$ as an ordered coprime factorization. For $x>0$ and $x<0$
one obtains, respectively,
\[
\begin{array}{ll}
x=aR^2,\ x-1=bS^2,&aR^2-bS^2=1,\\
x=-aR^2,\ x-1=-bS^2,&bS^2-aR^2=1,
\end{array}
\]
together with $a^2R^4-dU^2=9$. There are 18 branches. Direct exhaustion
of $(R,S,U)$ modulo the listed modulus eliminates 15:
\[
\begin{array}{c|c|c}
d&(a,b),\operatorname{sign}(x)&\text{modulus}\\ \hline
1&(1,1),+;\ (1,1),-&16;\ 8\\
2&(1,2),-;\ (2,1),+;\ (2,1),-&9;\ 8;\ 8\\
3&(1,3),+;\ (1,3),-;\ (3,1),+&9;\ 8;\ 8\\
6&(1,6),+;\ (1,6),-;\ (2,3),+;\ (2,3),-&9;\ 8;\ 8;\ 8\\
 &(3,2),-;\ (6,1),+;\ (6,1),-&8;\ 8;\ 8.
\end{array}
\]
The certificate exhausts every residue triple, making each row a finite
proof.

For $d=2,(a,b)=(1,2),x>0$,
\[
 R^2-2S^2=1,\quad R^4-2U^2=9,\quad
 (U-RS)(U+RS)=(R^2-9)/2.
\]
$R=1$ is impossible, $R=3$ gives $(S,U)=(2,6)$, and for $R\ge5$ both
factors are positive while $S>R/2$ makes the second factor larger than
their product.

For $d=3,(a,b)=(3,1),x<0$,
\[
 S^2-3R^2=1,\quad U^2=3(R^4-1),\quad
 (RS-U)(RS+U)=R^2+3.
\]
$R=0,2$ are impossible, $R=1$ gives $(2,0)$, and for $R\ge3$,
$RS>\sqrt3R^2>R^2+3$.

For $d=6,(a,b)=(3,2),x>0$,
\[
 3R^2-2S^2=1,\quad3R^4-2U^2=3,\quad
 (U-RS)(U+RS)=(R^2-3)/2.
\]
$R$ is odd; $R=1$ gives $(1,0)$, and the same factor-size contradiction
holds for $R\ge3$. These branches yield $t=6,y=\pm72$ and the separated
branch points $t=-6,0$. Adding all four roots proves the theorem.
\end{proof}

\section{Same-parameter elimination}
Mask 77 has support $\{0,2,3,6\}$ and mask 89 has support
$\{0,3,4,6\}$. The integral reflection $t=-u-6$ gives
\[
 t(t+3)(t+4)(t+6)=u(u+2)(u+3)(u+6).
\]
Thus the nonbranch parameter sets are $\{6\}$ and $\{-12\}$.

\begin{theorem}[Finite classification]\label{thm:54}
Among the 98 patterns left above, 26 contain mask 77, 25 contain mask 89,
and 7 contain both. All 44 patterns in the union fail the simultaneous
15-character conditions. Exactly 54 candidate patterns remain.
\end{theorem}
\begin{proof}
The union has $26+25-7=44$ elements. For a pattern containing only mask
77, substitute $t=6$ into all 15 character products; at least one other
product is nonsquare. For mask 89 use $t=-12$. If both occur, the
singleton candidate sets are disjoint. The certificate records every
failed mask, every common candidate set, and all 54 surviving identifiers.
It independently asserts 15 distinct masks in every input row.
\end{proof}

\section{A fourth-round gate: mask 102}

The remaining 54 patterns were ranked by two exact quantities: the number
of patterns containing a given four-root character, and whether its four
linear factors admit a pairing into two monic quadratics with constant
difference. Mask 102, supported at $\{1,2,5,6\}$, ranks first: it occurs
in 19 patterns and has
\[
 A=(t+1)(t+6),\qquad B=(t+2)(t+5),\qquad B-A=4.
\]

\begin{proposition}\label{prop:102}
The integral points on
\[
 y^2=(t+1)(t+2)(t+5)(t+6)
\]
are exactly
\[
 (-6,0),\quad(-5,0),\quad(-2,0),\quad(-1,0).
\]
\end{proposition}
\begin{proof}
For $t\leq-7$ or $t\geq0$, both $A$ and $B$ are positive.
Moreover $\gcd(A,B)\mid4$. If $AB$ is a square, the common positive
squarefree kernel is $d=1$ or $2$, so $A=dU^2,B=dV^2$.
For $d=1$,
\[
 (V-U)(V+U)=V^2-U^2=4.
\]
The only positive factor pair of equal parity is $(2,2)$, forcing $U=0$
and hence a branch point. For $d=2$ one obtains
$V^2-U^2=2$, impossible modulo 4. The six middle integers
$-6\leq t\leq-1$ give right sides $0,0,12,12,0,0$, respectively.
This proves the list.
\end{proof}

\begin{corollary}\label{cor:35}
Mask 102 eliminates 19 of the 54 patterns in
Theorem~\ref{thm:54}. Hence 35 candidate patterns remain after this
fourth-round gate.
\end{corollary}
\begin{proof}
All integral points of Proposition~\ref{prop:102} have
$t\in\{-6,\ldots,0\}$ and are therefore degenerate for the original
seven-term problem. The exact occurrence table contains mask 102 in
19 of the 54 rows. Thus no same-parameter candidate survives in those
rows.
\end{proof}

\section{A fifth-round gate: mask 108}
Mask 108 has support $\{2,3,5,6\}$ and occurs in 12 of the 35 patterns.

\begin{proposition}\label{prop:108}
The integral points on
\[
 y^2=(t+2)(t+3)(t+5)(t+6)
\]
are exactly
\[
 (-6,0),(-5,0),(-4,-2),(-4,2),(-3,0),(-2,0).
\]
\end{proposition}
\begin{proof}
Put $A=(t+2)(t+6)$ and $B=(t+3)(t+5)=A+3$.
For $t\leq-7$ or $t\geq-1$, both are positive and
$\gcd(A,B)\mid3$.  If $AB$ is a square, their common positive squarefree
kernel is $d=1$ or $3$, so $A=dU^2,B=dV^2$.
For $d=1$, $(V-U)(V+U)=3$ gives $(U,V)=(1,2)$, while $A=1$ would give
$(t+4)^2=5$, impossible modulo 8.  For $d=3$,
$(V-U)(V+U)=1$ forces $U=0$, contradicting $A>0$.
At the five middle integers $-6\leq t\leq-2$, direct substitution gives
the displayed six points.
\end{proof}

\begin{theorem}[Summary classification]\label{thm:summary}
Let $t\in\mathbf Z\setminus\{-6,-5,\ldots,0\}$.  If the seven nonzero
classes $[t],\ldots,[t+6]$ have affine rank at most two, then their
equality partition, modulo relabelling and reversal, is one of the 23
patterns in Table~\ref{tab:23}.  In particular, the exact reduction chain is
\[
651\longrightarrow343\longrightarrow284\longrightarrow98
\longrightarrow54\longrightarrow35\longrightarrow23.
\]
This theorem neither asserts that any of the 23 patterns is realizable nor
decides $R_2(7)$.
\end{theorem}
\begin{proof}
Affine rank zero or one is excluded by Lemma~\ref{lem:R16}, after applying
the common rational scaling, to an endpoint window.  Thus the actual affine
rank is two.  The map from the affinely spanning labels to squareclasses
therefore has rank two and is injective, so the labels are precisely the
actual squareclass equality blocks.  Proposition~\ref{prop:count} and
Lemma~\ref{lem:15} give 284 self-contained packets.  Lemma~\ref{lem:four}
leaves 98; Theorem~\ref{thm:54} leaves 54;
Proposition~\ref{prop:102} leaves 35; and Proposition~\ref{prop:108}, whose
points all have a zero among the seven original terms, eliminates 12 more.
The exact occurrence audit gives Table~\ref{tab:23}.
\end{proof}

\begin{table}[ht]
\centering
\caption{Final canonical candidates.  ID is the zero-based index in the
supplement's \texttt{unresolved\_patterns\_ranked}, ordered by its documented
difficulty score and then by restricted-growth word.}\label{tab:23}
\begin{tabular}{r@{\quad}l r@{\quad}l r@{\quad}l}
\toprule
ID&partition&ID&partition&ID&partition\\ \midrule
9&0012011&12&0012202&26&0112001\\
31&0001202&33&0001220&43&0100021\\
50&0111201&59&0012231&83&0112320\\
134&0012131&188&0112132&210&0121132\\
212&0121310&214&0122213&230&0012102\\
251&0102221&257&0112102&266&0120210\\
268&0121102&271&0122201&276&0010203\\
281&0102003&283&0111231&&\\ \bottomrule
\end{tabular}
\end{table}

\section{Supplement and reproducibility}
The accompanying supplement is identified by release
\texttt{paper-square-supplement-v0.5.0}.  Its root manifest is
\texttt{PAPER\_SQUARE\_SUPPLEMENT\_MANIFEST.json}; it records SHA-256 and
byte length for every generator, input/output certificate and regression
test, together with Python/SymPy versions and clean reproduction commands.
The manifest SHA-256 is
\texttt{6004afc5334bbea62969640079cad80764d2938d73b6fd62a85abc2249794588}.
Its release status is explicitly
\texttt{LOCAL\_RELEASE\_CANDIDATE\_NOT\_PUBLICLY\_ARCHIVED}; a journal
submission must replace that status with a frozen archive URL/DOI without
changing the hashed payload.

The frozen journal-candidate package is separately described by a
non-self-referential root manifest.  The file is
\begin{center}
\texttt{PAPER\_SQUARE\_SUBMISSION\_MANIFEST.json}
\end{center}
and its release identifier is \texttt{paper-square-submission-v0.6.4}.
It binds this source, the
bibliography, one reference PDF, the supplement manifest, and all submission
prose.  The reference PDF is byte-bound as an archived artifact, but PDF
metadata timestamps mean that an independent rebuild is required to agree in
mathematical content and diagnostics, not necessarily byte for byte.

The full regression command is stored verbatim in the manifest.  The finite
congruence certificate stores, for each of the 15 excluded mask-77 branches,
the modulus and exhaustive residue predicate.  The remaining three branches
carry symbolic factor identities and thresholds.

\paragraph{Online Resource 1 (supplementary code and certificates).}
The supplementary ZIP contains the exact Python generators, regression tests,
JSON certificates and nested manifest used to audit the finite pattern counts
and the mask-77, mask-89, mask-102 and mask-108 exclusions.  It contains no
additional mathematical claim beyond the theorem and claim boundary stated in
this article.  Before submission, the file metadata and upload form must add
the approved article title, journal name, author names, affiliations and
corresponding-author e-mail address.

\section{Prior-art boundary and limitations}
Work over number fields \cite{XarlesNF,GJXFive,BremnerSiksek,GJ2026,GJTho2026}
asks for terms that are themselves squares in a prescribed extension.
Our affine formulation allows Proposition~\ref{prop:scaling}'s common
rational scaling. Nearby algorithmic work concerns $S$-integral points
\cite{PethoEtAl}, simultaneous Pell equations \cite{MasserRickert},
and integral points in progression on congruent-number curves \cite{BST}.
Exact-quartic, $S$-integral-transform and Pell-system searches found no
published statement identical to Theorems~\ref{thm:H} and~\ref{thm:54}.
This is a \emph{not found} report, not proof of novelty.

The manifest and nested certificates bind SAFE and Round-04 inputs by
SHA-256 and store the
18 branches, structured factor identities, 44 same-parameter audits, and
54 remaining identifiers. The fourth-round certificate separately stores
the mask-102 ranking, proof data, the 19 excluded identifiers and the
35 intermediate survivors; the mask-108 certificate stores its complete
point proof data, 12 exclusions and 23 final candidates. An end-to-end test regenerates
651--343--284 and all $284\cdot15=4260$ characters.
We do not prove that the final 23 patterns are realized, or that no
seven-term rank-two progression exists. The 2026 papers have additional
rank and class-number hypotheses; none is inferred from a label here.

\section*{Statements and Declarations}

\paragraph{Funding.}
\textbf{AUTHOR CONFIRMATION REQUIRED.}  The human authors must insert and
approve a complete and accurate funding statement before submission; no
funding status is inferred here.

\paragraph{Competing interests.}
\textbf{AUTHOR CONFIRMATION REQUIRED.}  The human authors must insert and
approve a complete statement of relevant financial and non-financial
interests before submission; no absence of interests is inferred here.

\paragraph{Author identity and contributions.}
\textbf{AUTHOR CONFIRMATION REQUIRED.}  No author identity, affiliation,
corresponding-author information or contribution allocation is inferred.
Every listed author must satisfy the journal's authorship criteria, approve the
final manuscript and accept accountability for the work.

\paragraph{Use of large language models and AI-assisted tools.}
Large language models were used in this project for mathematical exploration,
generation and revision of computer code, literature-search assistance, and
generation and revision of manuscript and submission text.  No AI system is
an author.  Before submission, the human authors must independently verify
every proof, computation, code path, citation and statement, confirm that this
disclosure accurately describes the use made, approve the final text and
accept full responsibility for its integrity.

\paragraph{Data and code availability.}
The code, certificates and manifests are currently only a local release
candidate and have not been deposited in a public repository.
\textbf{AUTHOR CONFIRMATION REQUIRED.}  Before submission, the human authors
must approve an immutable public deposit, insert its genuine persistent
identifier here, and verify the downloaded archive against the published
SHA-256 manifest.  No DOI or public URL is claimed in this draft.

\paragraph{Ethics approval.}
Not applicable: this mathematical study involved no human participants,
animals or personal data.

\begingroup
\small
\bibliographystyle{plain}
\bibliography{references}
\endgroup
\end{document}
